\chapter{Заключение}
\label{ch:zaklyuchenie}

\section{Обобщение}
\label{sec:obobshtenie}

Настоящата работа постави и предостави опит за решаване на
задачата за \emph{предсказване на динамиката} на
партии на ударни инструменти в MIDI запис —
в частност, даване на предположение на силата (\eng{velocity}) на всяка нота
единствено от
структурни и времеви входни данни, при строга забрана за използване на съседни
стойности за сила. Задачата бе формализирана като изобразяване на структурно-времеви
контекст в нотна динамика (\eng{velocity}), с детерминистичен и вероятностен вариант.

Методологично работата се отличава от съществуващите разработки по това,
че третира времето като
условен вход, а не като изход, и затова вместо мрежовото представяне на нотите
използва редица в комбинация с непрекъснатата метрична фаза, както и други
входни характеристики, избягвайки артефактите от квантуване върху дуплетна мрежа на
триолни и неравноделни усещания. Дефинициите на характеристиките (групиране на
гласовете, толеранс за едновременност) бяха изведени управлявано от данните чрез
предварителен анализ на обучаващото подмножество.

Експериментално беше изградено и оценено семейство от модели. Табличният
LightGBM (градиентен бустинг с дървета) остава най-силен по \emph{точкова}
точност (MAE $18{,}02$) и е
интерпретируем, но свива динамичния обхват. Неавторегресивният
трансформър побеждава наивните базови модели, и макар да отстъпва на дърветата по
MAE, запазва значително по-добре динамичния обхват. Решаващата стъпка бяха
\emph{вероятностните глави} върху същия трансформър ,,гръбнак'':
категорийният изход съчетава най-добрата точкова точност сред главите и — чрез семплиране —
почти идеално съответствие с истинското разпределение на търсената сила в тестовото подмножество
(histI $0{,}947$), което е същинската цел на задачата и е недостижимо за
точков модел.

Анализът на интерпретируемостта показа, че моделът е научил реални барабанни
явления: динамична идентичност на инструментите, метричната йерархия, ритмични конвенции по стил
и — при семплиране — разпределението на призрачните ноти и акцентите. Беше количествено
демонстрирана регресията към средното при точковата оценка и нейното преодоляване
чрез семплиране. Проверката за умението за генерализация чрез обучаване по
подмножество, нарочно пропускащо определени барабанисти, показа, че
относителната форма на динамичния обхват се пренася към нечути изпълнители, но абсолютното
им динамично ниво — не.

\section{Насоки за бъдещи подобрения}
\label{sec:budeshti}

Работата очертава няколко естествени продължения:
\begin{itemize}
  \item \textbf{Условност по абсолютно ниво:} Тъй като абсолютната сила на
        удара е индивидуална черта, подаването на целево ниво на силата за
        изпълнението (или няколко калибриращи ноти), или пък предсказването на
        \emph{нормализирана в рамките на записа} динамика с последващо
        премащабиране, би премахнало основния източник на грешка при нечуван
        изпълнител.
  \item \textbf{Преизграждане върху една бленда:} Систематичните погласови
        отмествания се дължат на записването през 43 различни бленди на комплекта
        електронни барабани, които по различен начин често пренасочват
        едни и същи падове към различни MIDI височини и внася шум във входната характеристика
        „глас''. Преизграждане върху една
        канонична бленда (напр.\ често срещаният „Acoustic Kit'') е чиста проверка дали тези
        отмествания изчезват и намалява излишъка в данните.
  \item \textbf{Използване на повече тренировъчни данни:} Записи от повече различни барабанисти,
        стилове, размери и фрази, или пък цели песни, биха помогнали за по-качествено обучение
        и допринесли за по-добра генерализация.
  \item \textbf{Съвместно моделиране на време и динамика:} Естествено разширение е
        моделът да предсказва едновременно микротайминг и динамика.
  \item \textbf{Емпирична оценка:} Формален A/B тест със слушатели би измерил
        възприеманото качество на предсказаната динамика отвъд числовите метрики.
\end{itemize}
